\SourceLeaf{232}%
gam determinatam eligimus, in qua figura concipi debet. In sphaera auxiliari semper plagam exteriorem, a centro aversam, adhibebimus: in superficie curva etiam plaga exterior sive quae tamquam exterior consideratur, adoptari potest, vel potius plaga eadem, a qua normalis erecta concipitur; manifesto enim respectu similitudinis figurarum nihil mutatur, si in superficie curva tum figura ad plagam oppositam transfertur, tum normalis, dummodo ipsius imago semper in eadem plaga superficiei sphaericae depingatur.

Signum positivum vel negativum, quod pro situ figurae infinite parvae \textit{mensurae} curvaturae adscribimus, etiam ad curvaturam integram figurae finitae in superficie curva extendimus. Attamen si argumentum omni generalitate amplecti suscipimus, quaedam dilucidationes requiruntur, quas hic breviter tantum attingemus. Quamdiu figura in superficie curva ita comparata est, ut singulis punctis intra ipsam puncta \textit{diversa} in superficie sphaerica respondeant, definitio ulteriore explicatione non indiget. Quoties autem conditio ista locum non habet, necesse erit, quasdam partes figurae in superficie sphaerica bis vel pluries in computum ducere, unde, pro situ simili vel opposito, vel accumulatio vel destructio oriri poterit. Simplicissimum erit in tali casu, figuram in superficie curva in partes tales divisam concipere, quae singulae per se spectatae conditioni illi satisfaciant, singulis tribuere curvaturam suam integram, quantitate per aream figurae in superficie sphaerica respondentis, signo per situm determinatis, ac denique figurae toti adscribere curvaturam integram ortam per additionem curvaturarum integrarum, quae singulis partibus respondent. Generaliter itaque curvatura integra figurae est \(=\int k\,\mathrm{d}\sigma\), denotante \(\mathrm{d}\sigma\) elementum areae figurae, \(k\) mensuram curvaturae in quovis puncto. Quod vero attinet ad repraesentationem geometricam huius integralis, praecipua huius rei momenta ad sequentia redeunt. Peripheriae figurae in superficie curva (sub restrictione art. 3) semper respondebit in superficie sphaerica linea in se ipsam rediens. Quae si se ipsam nullibi intersecat, totam superficiem sphaericam in duas partes dirimet, quarum altera respondebit figurae in superficie curva, et cuius area, positive vel negative accipienda, prout respectu peripheriae suae similiter iacet ut figura in superficie curva respectu suae, vel inverse, exhibebit posterioris curvaturam integram. Quoties vero linea ista se ipsam semel vel pluries secat, exhibebit figuram complicatam, cui tamen area certa aeque legitime tribui potest, ac figuris absque nodis, haecque area, rite intellecta, semper valorem iustum curvaturae integrae exhibebit. Attamen uberio%
